\documentclass[11pt, a4paper]{article}

\usepackage[T1]{fontenc}
\usepackage[utf8]{inputenc}
\usepackage{tgpagella}
\usepackage[spanish, es-noshorthands]{babel}
\usepackage{amsmath, amssymb, bm}
\usepackage{tabularx}
\usepackage{booktabs}
\usepackage[most]{tcolorbox}
\usepackage[margin=2.5cm]{geometry}
\usepackage{ragged2e}
\usepackage{fancyhdr}
\usepackage[american]{circuitikz}

\pagestyle{fancy}
\fancyhf{}
\setlength{\headheight}{16pt}
\lhead{\textbf{UCB Sede Tarija} | Ing. Mecatrónica}
\chead{}
\rhead{IMT-121 Circuitos Electrónicos I | Análisis RC Guiado}
\lfoot{Prof: M.Sc. Ing. Bernardo Quiroga Turdera}
\rfoot{Página \thepage}
\renewcommand{\headrulewidth}{0.4pt}

\definecolor{ucbblue}{RGB}{0, 51, 102}

\begin{document}

\begin{tcolorbox}[colback=ucbblue!5!white, colframe=ucbblue, title=\textbf{Hoja de Análisis Guiado: Transitorios de Primer Orden RC}]
    \centering \textbf{Sistematización de Cálculos y Gráficas}
\end{tcolorbox}

\noindent\textbf{Instrucciones:} Dado un circuito serie RC genérico sometido a una conmutación en $t=0$, proceda paso a paso para construir y validar la respuesta del sistema.

\vspace{0.4cm}
\section*{Parte 1: Análisis Dinámico y Continuidad}
\begin{enumerate}
    \item Basado en el análisis de la red en estado estacionario pre-conmutación ($t<0$), determine el voltaje inicial del capacitor:\\
    $v_C(0^-) =$ \rule{4cm}{0.4pt}
    
    \item Aplicando el principio físico de continuidad (el capacitor se opone a cambios bruscos de tensión por su inercia de campo eléctrico):\\
    $v_C(0^+) =$ \rule{4cm}{0.4pt}
\end{enumerate}

\section*{Parte 2: Análisis Estacionario Final}
\begin{enumerate}
    \item Conecte la red post-conmutación. Asuma que transcurrió mucho tiempo ($t \to \infty$) y el capacitor está completamente adaptado. En ese caso, la derivada temporal $\frac{dv_C}{dt} = $ \rule{2cm}{0.4pt}.
    \item Por tanto, el capacitor se modela topológicamente como un \rule{4cm}{0.4pt}.
    \item El voltaje de estado estable en el capacitor será:\\
    $v_C(\infty) =$ \rule{4cm}{0.4pt}
\end{enumerate}

\section*{Parte 3: Modelo Dinámico e Inercia}
\begin{enumerate}
    \item Apague todas las fuentes independientes de la red post-conmutación. Determine la resistencia equivalente vista desde los bornes del capacitor:\\
    $R_{eq} =$ \rule{4cm}{0.4pt}
    
    \item Calcule la constante de tiempo:\\
    $\tau = R_{eq}C =$ \rule{4cm}{0.4pt} (segundos)
\end{enumerate}

\section*{Parte 4: Construcción y Verificación de Respuesta}
Sustituya en la forma canónica: $v_C(t) = v_C(\infty) + [v_C(0^+) - v_C(\infty)]e^{-t/\tau}$
\begin{equation}
    v_C(t) = \rule{8cm}{0.4pt}
\end{equation}

\textbf{Comprobación de Extremos (Sanity Check):}
\begin{itemize}
    \item Evaluar en $t=0$: ¿Coincide con $v_C(0^+)$? \rule{3cm}{0.4pt}
    \item Evaluar en $t \to \infty$: ¿Coincide con $v_C(\infty)$? \rule{3cm}{0.4pt}
\end{itemize}

\textbf{Interpretación Física:}
\begin{itemize}
    \item El término $\frac{dv_C}{dt}$ en la EDO representa \rule{8cm}{0.4pt}.
    \item Conforme la ecuación, el voltaje del capacitor \textbf{[Sube / Baja]} progresivamente desde \rule{2cm}{0.4pt} V hacia \rule{2cm}{0.4pt} V.
\end{itemize}

\end{document}
